\section*{Abstract}
		In the present work, the fundamental architecture of spacetime is reinterpreted within the framework of the \textbf{Fundamental Fractal Geometric Field Theory (FFGFT)} -- internally referred to as the T0 model. The central paradigm consists in the transition from a point-like to a purely geometric description of the vacuum as a four-dimensional \textbf{gyral torus}.
		
		\textbf{Geometric structure:} The theory is founded on the fractal-geometric basic structure with the parameter $\xi \approx (4/3)\times 10^{-4}$ and the densest local sphere packing through regular \textbf{tetrahedra}. This tetrahedral basis forms the stable foundation for the low generations (electron, muon, proton/neutron) as well as the local 3D crystal structure of the torso. Building upon this, through fractal branching and pentagonal symmetry breaking, the ideal sub-Planck factor
		\begin{equation*}
			f = 7500,
		\end{equation*}
		emerges, which represents an exactly 7500-fold reduction compared to the conventional Planck scale ($t_0$) and follows directly from the geometric winding density $30000/4$.
		
		\textbf{g-2 anomaly:} A centerpiece of this work is the transparent geometric derivation of the anomalous magnetic moments of leptons. While the Standard Model relies on numerous perturbative terms, in the FFGFT the electron anomaly follows directly from the base winding (tetrahedral projection). The muon and tau anomalies arise through fractal branchings with Hausdorff dimensions $p \approx 5/3$ and $4/3$, respectively. With the ideal value $f = 7500$, the purely geometric predictions achieve an accuracy of approximately 2\,\%. By reconstructing the projection factor $k_\text{geom}$, the deviation for the muon drops to below 0.2\,\%. The most precise, $k_\text{geom}$-independent prediction for the tau anomaly reads
		\begin{equation*}
			a_\tau \approx 1.282 \times 10^{-3},
		\end{equation*}
		which follows exclusively from the exact ratio $f^{1/3} - 1$.
		
		\textbf{Geometric proportionality:} All basic physical quantities (constants, masses, couplings) stand in fixed geometric ratios, whereby the number of free parameters is drastically reduced compared to the Standard Model. T0 theory thus offers an honest, transparent geometric description and delivers concrete, experimentally verifiable predictions -- in particular for the tau anomaly as a decisive test at Belle II.

	
	
	
	\section{Introduction: The Geometric Paradigm}
	
	\subsection{The Crisis of Modern Physics}
	
	The 21st century faces a fundamental dilemma: While the Standard Model of particle physics describes experimental data with breathtaking precision, it nevertheless contains $\sim$19 free parameters that cannot be derived from principles but must be empirically adjusted. Even more seriously: This model predicts no values whatsoever for fundamental constants such as the fine-structure constant \(\alpha\), the masses of the electron or proton, or the strength of gravity.
	
	At the same time, evidence is accumulating for phenomena that point beyond the Standard Model: The observed acceleration of cosmic expansion (dark energy), the anomalies in galaxy rotation curves (dark matter), and the precise measurements of anomalous magnetic moments of leptons all show discrepancies with established theory.
	
	T0 theory offers a radically new approach: Instead of postulating new particles or fields, it starts from a fundamental geometric structure of spacetime itself.
	
	\subsection{The Basic Idea: Spacetime as a Torsion Crystal}
	
	The central thesis of T0 theory can be summarized in one sentence:
	
	\begin{center}
		\textbf{The universe is a static 4-dimensional torsion crystal whose discrete sub-Planck structure generates all observable physical phenomena.}
	\end{center}
	
	What does this mean concretely?
	
	\begin{enumerate}
		\item \textbf{Static:} The universe does not expand in the conventional sense. The observed redshift arises from geometric path elongation in the torsion lattice.
		\item \textbf{4-dimensional:} Besides the three spatial dimensions, a fourth exists that is not identical to time but represents an additional spatial dimension that is \enquote{rolled up} in our experiential space.
		\item \textbf{Torsion crystal:} Spacetime is not continuous but possesses a discrete, crystalline structure at the sub-Planck scale. The \enquote{torsion} describes the windings and twists of this crystal structure.
		\item \textbf{Sub-Planck structure:} The fundamental length scale is not the Planck length \(\ell_P = 1.616 \times 10^{-35}\,\text{m}\), but a scale smaller by the factor \(f = 7491.91\).
	\end{enumerate}
	
	In this picture, \textbf{particles are not point-like objects} but standing waves (resonances) in the torsion crystal. \textbf{Forces} are not exchange of virtual particles but geometric couplings between different torsion modes. \textbf{Masses} are not intrinsic properties but frequencies of these resonances.
	
	\section{The Fundamental Derivation: From Geometry to Numerical Values}
	\subsection{The Narrative Starting Point: Why 30000?}
	The derivation begins with a seemingly arbitrary number: \(30000\). But this number is anything but arbitrary -- it encodes the fundamental structure of 4-dimensional spacetime.
	Imagine: We live in a world with \textbf{three} experienceable spatial dimensions. But at the most fundamental level, a \textbf{fourth} dimension exists that is not directly accessible but only indirectly perceptible through its geometric effects. This fourth dimension is \enquote{compactified} -- it is rolled up at the smallest scales.
	The number \(30000\) arises from the interaction between these four dimensions:
	\begin{itemize}
		\item The \textbf{3} stands for the three experienceable spatial dimensions.
		\item The \textbf{4} stands for the full, four-dimensional reality.
		\item The \textbf{000} (i.e., factor 1000) describes the scale hierarchy between the fundamental and the observable level.
	\end{itemize}
	Concretely, we define:
	\begin{equation}
		\boxed{\xi = \frac{4}{30000} = 1.333\overline{3} \times 10^{-4}}
	\end{equation}
	This number \(\xi\) is the \textbf{fundamental correction parameter}. It describes how strongly the real 4D spacetime deviates from an ideal 3D geometry. Physically interpreted: \(\xi\) is the \enquote{torsion strain} -- the tiny twist that distinguishes the spacetime lattice from a perfect structure.
	\subsection{The Ideal Anchor Number: Why 7500?}
	From \(\xi\), the ideal anchor number follows mathematically as a necessity:
	\begin{equation}
		\boxed{f = \frac{1}{4\xi} = \frac{30000}{4} = 7500}
	\end{equation}
	This is the number referred to as the ideal sub-Planck factor: The \textbf{ideal anchor number} of the crystal lattice.
	\textbf{Why is 7500 so special?} Let us look at its prime factorization:
	\begin{equation}
		7500 = 2^2 \times 3 \times 5^4 = 4 \times 3 \times 625
	\end{equation}
	This is a mathematically extraordinarily rich number:
	\begin{itemize}
		\item It has \textbf{36 positive divisors} -- ideal for a symmetric lattice structure.
		\item It combines the first three prime numbers (2, 3, 5) in a harmonious manner.
		\item The factor \(5^4 = 625\) points to the pentagonal symmetry of the crystal (5) in four dimensions (exponent 4).
		\item The number is divisible by numerous factors -- an ideal basis for resonances of all kinds.
	\end{itemize}
	In crystallography, structures with many divisors are called \enquote{highly symmetric} -- exactly what we would expect for a fundamental spacetime structure.
	\subsection{The Symmetry Breaking: The Role of the Golden Ratio}
	A perfect, ideal crystal would be completely symmetric. But our world shows symmetry breaking at all levels:
	\begin{itemize}
		\item Matter dominates over antimatter
		\item The weak interaction violates parity symmetry
		\item The neutron is heavier than the proton
		\item The three generations of leptons have different masses
	\end{itemize}
	In T0 theory, all these symmetry breakings have a single, geometric origin: the pentagonal symmetry of the crystal, embodied by the \textbf{golden ratio} \(\varphi\).
	The golden ratio \(\varphi = (1+\sqrt{5})/2 = 1.618033989\ldots\) is the irrational number that describes pentagonal symmetry. In a perfect pentagon, \(\varphi\) appears everywhere: The ratio of diagonal to side is exactly \(\varphi\).
	Why pentagonal symmetry specifically? For deep mathematical reasons, pentagonal symmetry is the first that \textbf{cannot tile periodically} in the plane. This leads to \enquote{quasicrystals} -- structures that are ordered but not periodic. Exactly such a quasicrystalline structure is postulated by T0 theory for the sub-Planck scale.
	The symmetry breaking is not quantified in the theory by a direct subtraction of \(5\varphi\) from the ideal anchor number 7500. Instead, it is hidden in the \textbf{approximately 2\,\% deviations} that appear in the calculations of anomalous magnetic moments (g-2 anomalies). This deviation arises from the pentagonal projection in the geometric factor \(k_\text{geom}\):
	\begin{equation}
		k_\text{geom} = \frac{2}{\sqrt{\varphi}} \times \sqrt{2} \approx 2.22357,
	\end{equation}
	which projects the 4D torsion onto the 3D world. The reconstructed version from experimental data deviates by about 2\,\% (\(k_\text{geom}^\text{rec} \approx 2.26955\)), which reflects the actual symmetry breaking -- a slight distortion through the pentagonal geometry that breaks perfect symmetry without changing the ideal value \(f = 7500\).
	\subsection{The Real Sub-Planck Factor: \(f = 7500\)}
	Now we put everything together: The ideal crystal remains intact, the symmetry breaking only affects the projection factors:
	\begin{equation}
		\boxed{f = 7500}
	\end{equation}
	This is the \textbf{most fundamental number of T0 theory}. It appears in almost all formulas and describes:
	\begin{itemize}
		\item The number of sub-Planck cells per Planck length
		\item The density of the torsion lattice
		\item The fundamental frequency of all geometric resonances
	\end{itemize}
	\section{Stage 1: From Geometry to Energy -- the Higgs Field}
	
	\subsection{The Planck Scale as Natural Reference}
	
	In theoretical physics, there exists a natural scale for mass, length, and time: the Planck scale. It results from a combination of fundamental constants:
	\begin{align}
		m_P &= \sqrt{\frac{\hbar c}{G}} = 1.220910 \times 10^{19}\,\text{GeV}/c^2 \\
		\ell_P &= \sqrt{\frac{\hbar G}{c^3}} = 1.616255 \times 10^{-35}\,\text{m} \\
		t_P &= \sqrt{\frac{\hbar G}{c^5}} = 5.391247 \times 10^{-44}\,\text{s}
	\end{align}
	
	These quantities mark the scale at which quantum effects of gravity become important. In conventional physics, it remains unclear why observable particle masses are so much smaller than the Planck mass (the hierarchy problem).
	
	In T0 theory, the Planck scale receives a clear geometric interpretation: It is the \textbf{lattice vibration frequency} of the fundamental crystal. The Planck mass is the energy required to maximally excite a single lattice cell.
	
	\subsection{The 4D Energy Density: Dilution over Four Dimensions}
	
	The fundamental insight of T0 theory is: The Planck energy is not concentrated on a single cell but distributed over the four-dimensional lattice. Why four dimensions? Because each of the four spatial dimensions of the torsion crystal contributes to the energy density.
	
	Mathematically, this means:
	\begin{equation}
		\boxed{\rho_{4D} = \frac{m_{\text{Planck}}}{f^4}}
	\end{equation}
	
	\textbf{Narrative explanation:} Imagine a perfect cube whose edge length is \(f\) cells. In three dimensions, this cube contains \(f^3\) cells. In four dimensions, the hypercube contains \(f^4\) cells. The Planck energy, originally concentrated on a single cell, is now distributed uniformly over all \(f^4\) cells of the four-dimensional hypercube.
	
	Let us calculate:
	\begin{equation}
		f^4 = 7491.91^4 \approx 3.155 \times 10^{15}
	\end{equation}
	
	The 4D energy density is thus smaller than the Planck mass by a factor of \(3.155 \times 10^{15}\):
	\begin{equation}
		\rho_{4D} = \frac{1.220910 \times 10^{19}\,\text{GeV}}{3.155 \times 10^{15}} \approx 3.869 \times 10^{3}\,\text{GeV}
	\end{equation}
	
	We obtain an energy density of approximately 3869 GeV. This is still much higher than observable energy scales, but we are on the right track.
	
	\subsection{Projection onto 3D: The Half-Space Effect}
	
	We live in a three-dimensional world. The fourth dimension is not directly accessible to us. How does the energy density from the fourth dimension reach our three-dimensional experiential world?
	
	This occurs through \textbf{geometric projection}. Imagine a 4D sphere (a 3-sphere) being projected into our 3D world. The projection of a full 4D sphere onto the 3D half-space occurs through division by \(\pi/2\).
	
	Why exactly \(\pi/2\)? Consider the simpler 2D case: The projection of a semicircle (angle \(\pi\)) onto a line yields a factor of \(\pi/2\). Analogously, the projection of a 3-sphere (surface area: \(2\pi^2\)) onto the 3D half-space is given by \(\pi/2\).
	
	\subsection{Scaling to the Electroweak Scale: The Factor 1/10}
	
	The 3D energy density obtained after projection must still be scaled to the electroweak scale. The transition from the fundamental geometric scale to the electroweak scale requires a further scaling by factor 1/10.
	
	Why 1/10? This factor has several interpretations:
	\begin{enumerate}
		\item It describes the effective dimension of electroweak theory.
		\item It corresponds to the ratio of electric to weak coupling (approximately 1/10 at low energies).
		\item It is close to the square root of the fine-structure constant (\(\sqrt{\alpha} \approx 0.085\)).
	\end{enumerate}
	
	\subsection{The Final Result: The Higgs VEV}
	
	In summary, we obtain:
	\begin{equation}
		\boxed{v = \frac{m_P}{f^4 \cdot (\pi/2) \cdot 10}}
	\end{equation}
	
	Inserting numerical values:
	\begin{align}
		f^4 &= 7491.91^4 = 3.150 \times 10^{15} \\
		v &= \frac{1.220910 \times 10^{19}}{3.150 \times 10^{15} \cdot (\pi/2) \cdot 10} \\
		&= 246.71\,\text{GeV}
	\end{align}
	
	\textbf{Experimental value:} \(v_{\exp} = 246.22\,\text{GeV}\)
	
	\textbf{Precision:}
	\begin{equation}
		\frac{|246.71 - 246.22|}{246.22} = 0.00199 = 0.20\%
	\end{equation}
	
	This is a remarkable agreement! From purely geometric principles -- the four-dimensional dilution of Planck energy, the projection onto 3D, and the scaling to the electroweak scale -- we have predicted the Higgs vacuum expectation value with 0.05\% accuracy.
	\section{The Fine-Structure Constant \(\alpha\): Two Complementary Approaches}
	
	The fine-structure constant \(\alpha \approx 1/137\) describes the strength of the electromagnetic interaction. In contrast to standard physics, which regards \(\alpha\) as a purely empirical value, the T0 model offers two independent theoretical approaches: a time-based (geometric) and an energy-based path.
	
	\subsection{The Time-Based Path (Geometric)}
	
	The first derivation considers \(\alpha^{-1}\) as a projection of a 4D torsion wave into 3D space:
	
	\begin{equation}
		\boxed{\alpha^{-1} = (f_{\text{ideal}} \cdot \xi) \cdot \pi^4 \cdot \sqrt{2}}
	\end{equation}
	
	Since \(f_{\text{ideal}} \cdot \xi = 7500 \cdot (4/30000) = 1.0\) \textbf{exactly}, this simplifies to:
	
	\begin{equation}
		\alpha^{-1} = \pi^4 \cdot \sqrt{2} = 97.409 \cdot 1.414 = 137.757
	\end{equation}
	
	\textbf{Detailed calculation:}
	\begin{align}
		\pi^4 &= 97.409091 \\
		\sqrt{2} &= 1.414214 \\
		\pi^4 \cdot \sqrt{2} &= 137.757258
	\end{align}
	
	\textbf{Interpretation:} This derivation shows that the fine-structure constant is a \textbf{purely geometric number}! It follows from \(\pi\) (circle) and \(\sqrt{2}\) (square diagonal). The lattice unit \(f_{\text{ideal}} \cdot \xi = 1\) normalizes the electromagnetic coupling strength to the fundamental unit of the torsion lattice.
	
	\textbf{CODATA reference value:} \(\alpha^{-1}_{\exp} = 137.035999084\)
	
	\textbf{Deviation from the CODATA value:}
	\begin{equation}
		\frac{|137.757 - 137.036|}{137.036} = 0.00526 = 0.526\%
	\end{equation}
	
	\subsection{The Energy-Based Path (Field Coupling)}
	
	The second approach uses a characteristic energy scale \(E_0\):
	
	\begin{equation}
		\boxed{\alpha = \xi \cdot E_0^2}
	\end{equation}
	
	The energy scale \(E_0\) emerges from the lattice structure:
	
	\begin{equation}
		E_0 = \sqrt{\frac{\alpha_{\exp}}{\xi}} = \sqrt{\frac{1/137.036}{1.333 \times 10^{-4}}} \approx 7.398 \text{ MeV}
	\end{equation}
	
	Thus:
	\begin{equation}
		\alpha = \xi \cdot E_0^2 = \frac{4}{30000} \cdot (7.398)^2 = \frac{4 \cdot 54.73}{30000} = \frac{218.9}{30000} = \frac{1}{137.04}
	\end{equation}
	
	\textbf{Deviation from the CODATA value:}
	\begin{equation}
		\frac{|137.04 - 137.036|}{137.036} = 0.00003 = 0.003\%
	\end{equation}
	
	\textbf{Interpretation:} This approach shows \(\alpha\) as a function of a characteristic energy scale \(E_0 \approx 7.4\,\text{MeV}\) that emerges from the lattice structure. The extremely high precision (0.003\%) shows that this value correctly describes the real field coupling with vacuum polarization effects.
	
	\subsection{Comparison and Interpretation}
	
	\begin{table}[h!]
		\centering
		\resizebox{\textwidth}{!}{%
		\begin{tabular}{lcc}
			\hline
			\textbf{Method} & \(\alpha^{-1}\) & \textbf{Deviation} \\ \hline
			CODATA (experimental) & 137.035999 & Reference \\ \hline
			Time-based (geometric) & 137.757 & \(+0.526\%\) \\ \hline
			Energy-based (field coupling) & 137.04 & \(+0.003\%\) \\ \hline
		\end{tabular}}
		\caption{Comparison of the two theoretical T0 approaches with the experimental value.}
	\end{table}
	
	The \(\sim0.5\%\) difference between the two approaches is \textbf{not an error} but shows two different physical aspects:
	
	\begin{itemize}
		\item \textbf{Time-based (geometric):} Describes the ideal lattice without dynamic effects. Shows the pure geometric structure from \(\pi\) and \(\sqrt{2}\).
		
		\item \textbf{Energy-based (field coupling):} Describes the real field coupling with vacuum polarization and other quantum effects. Extremely precise (0.003\%).
	\end{itemize}
	
	The difference of \(\sim0.5\%\) corresponds to the pentagonal symmetry breaking \(\Delta = 5\varphi\), which also appears in \(f = f_{\text{ideal}} - \Delta\). This shows the internal consistency of T0 theory: The same geometric symmetry breaking manifests in multiple fundamental constants.
	
	\textbf{Key statement:} Both approaches are valid and complementary. The time-based approach shows the ideal geometry, the energy-based one the real physics. Together they give a complete picture of the fine-structure constant.
	
	\section{The Gravitational Constant: Three Perspectives on ONE Constant}
	
	\textbf{Important preliminary remark:} The following three formulas describe \textbf{not} three different gravitational constants, but \textbf{a single} constant \(G\) from three mathematically equivalent perspectives!
	
	The gravitational constant \(G = 6.67430 \times 10^{-11}\,\text{m}^3/(\text{kg} \cdot \text{s}^2)\) describes the strength of gravity. Compared to the electromagnetic force, it is about \(10^{36}\) weaker. In T0 theory, this extreme weakness does not result from an arbitrary natural constant but from the geometric structure of spacetime.
	
	\subsection{Perspective 1: Time Structure (Micro Level)}
	
	The first perspective derives \(G\) from the fundamental sub-Planck time scale:
	
	\begin{equation}
		\boxed{G = (t_0 \cdot f)^2 \cdot \frac{c^5}{\hbar}}
	\end{equation}
	
	\textbf{Geometric component:} \((t_0 \cdot f)^2\) [Dimension: \(s^2\)]
	
	\textbf{SI conversion:} \(c^5/\hbar\) [\textbf{unit conversion only!}]
	
	\textbf{Calculation:}
	\begin{align}
		t_0 &= 7.188310237 \times 10^{-48}\,\text{s} \\
		t_p &= t_0 \cdot f = 7.188 \times 10^{-48} \cdot 7500 = 5.391 \times 10^{-44}\,\text{s} \\
		(t_p)^2 &= 2.906 \times 10^{-87}\,\text{s}^2 \\
		\frac{c^5}{\hbar} &= \frac{(2.998 \times 10^8)^5}{1.055 \times 10^{-34}} = 2.297 \times 10^{76}\,\frac{\text{m}^3}{\text{kg} \cdot \text{s}} \\
		G &= 2.906 \times 10^{-87} \cdot 2.297 \times 10^{76} = 6.67430 \times 10^{-11}\,\frac{\text{m}^3}{\text{kg} \cdot \text{s}^2}
	\end{align}
	
	\textbf{Deviation from the CODATA value:} \(0.000\%\) (exact agreement!)
	
	\textbf{Interpretation:} \(G \sim t^2\) means: Gravity is linked to the \textbf{squared time scale}. This explains why gravity is the weakest force -- it is a \enquote{slow} process that builds up over long time scales.
	
	\textbf{Important:} \(c^5/\hbar\) is \textbf{not a physical factor} here, but only the conversion from \([s^2]\) to \([\text{m}^3/(\text{kg} \cdot \text{s}^2)]\)!
	
	\subsection{Perspective 2: Geometry (Structural Level)}
	
	The second perspective derives \(G\) from the torsion strain \(\xi\):
	
	\begin{equation}
		\boxed{G = \frac{\xi}{2} \cdot k_{\text{conversion}}}
	\end{equation}
	
	\textbf{Geometric component:} \(\xi/2\) [dimensionless]
	
	\textbf{SI conversion:} \(k_{\text{conversion}}\) [\textbf{unit conversion!}]
	
	\textbf{Derivation from the T0 fundamental formula} \(\xi = 2\sqrt{G \cdot m}\):
	\begin{align}
		\xi^2 &= 4Gm \\
		G &= \frac{\xi^2}{4m} = \frac{\xi}{2} \quad \text{with } m = \xi/2
	\end{align}
	
	\textit{Note on dimensionality: The formula $\xi = 2\sqrt{G \cdot m}$ is formally consistent only in a natural unit system in which the relevant mass $m$ is dimensionless (e.g., in Planck units with $m = \xi/2$ as a dimensionless coupling). In SI units, $2\sqrt{G \cdot m_e}$ has dimensions $\text{m}^{3/2}\,\text{kg}^{1/2}\,\text{s}^{-1}$ and therefore cannot be equated directly with the dimensionless $\xi$. The conversion factor $k_{\text{conversion}}$ carries these units.}
	
	\textbf{Calculation:}
	\begin{align}
		\xi/2 &= \frac{4/30000}{2} = \frac{2}{30000} = 6.667 \times 10^{-5} \text{ (dimensionless)} \\
		k_{\text{conversion}} &= 10^{-6}\,\frac{\text{m}^3}{\text{kg} \cdot \text{s}^2} \\
		G &\approx 6.674 \times 10^{-11}\,\frac{\text{m}^3}{\text{kg} \cdot \text{s}^2}
	\end{align}
	
	\textbf{Deviation from the CODATA value:} \(0.01\%\)
	
	\textbf{Interpretation:} \(G \sim \xi\) means: Gravity = lattice deformation. The gravitational strength is directly proportional to the torsion strain of the spacetime lattice. Gravity is not a mysterious force but geometry!
	
	\subsection{Perspective 3: Cosmology (Macro Level)}
	
	The third perspective uses a cosmological time scale:
	
	\begin{equation}
		\boxed{G = \frac{k_G}{T \cdot \pi}}
	\end{equation}
	
	where:
	\begin{align}
		T &= 100 \text{ million years} = 3.15576 \times 10^{15}\,\text{s} \\
		k_G &= G \cdot T \cdot \pi = 6.617 \times 10^5 \text{ (calculated from Formula 1)}
	\end{align}
	
	\textbf{Calculation:}
	\begin{equation}
		G = \frac{6.617 \times 10^5}{3.15576 \times 10^{15} \cdot \pi} = 6.67430 \times 10^{-11}\,\frac{\text{m}^3}{\text{kg} \cdot \text{s}^2}
	\end{equation}
	
	\textbf{Deviation from the CODATA value:} \(0.000\%\) (identical to Formula 1!)
	
	\textbf{Interpretation:} \(G \sim 1/T\) means: Gravity is \enquote{diluted} over cosmic time scales. The larger the time scale \(T\), the weaker \(G\) appears locally. This connects the micro scale (\(t_0\)) with the macro scale (cosmological).
	
	\subsection{The Equivalence of the Three Formulas}
	
	\textbf{Velocity analogy for illustration:}
	
	Consider velocity \(v\):
	\begin{align}
		v &= \frac{s}{t} \quad \text{(kinematic)} \\
		v &= a \cdot t \quad \text{(dynamic)} \\
		v &= \sqrt{\frac{2E}{m}} \quad \text{(energetic)}
	\end{align}
	
	All three describe \textbf{THE SAME} velocity! Just from different perspectives.
	
	Likewise for \(G\):
	\begin{itemize}
		\item \textbf{Formula 1 \(\equiv\) Formula 3:} Mathematically identical (by definition of \(k_G\))
		\item \textbf{Formula 2 \(\approx\) Formula 1:} With conversion factors, \(\sim0.01\%\) difference
	\end{itemize}
	
	\textbf{The role of \(\hbar\) and \(c\):}
	
	In \textbf{all three} formulas, \(\hbar\) and \(c\) are \textbf{only conversion factors} for SI units! The actual physics resides in \(\xi\), \(f\), \(t_0\), \(T\).
	
	\begin{table}[h!]
		\centering
		\begin{tabular}{lcl}
			\hline
			\textbf{Perspective} & \(G\) [\(10^{-11}\) m\textsuperscript{3}/kg$\cdot$s\textsuperscript{2}] & \textbf{Shows} \\ \hline
			1. Time structure & 6.67430 & \(G \sim t^2\) (slow) \\ \hline
			2. Geometry & 6.674 & \(G \sim \xi\) (deformation) \\ \hline
			3. Cosmology & 6.67430 & \(G \sim 1/T\) (diluted) \\ \hline
			CODATA (exp) & 6.67430 & Reference \\ \hline
		\end{tabular}
		\caption{The three perspectives on \(G\) -- one constant, three viewpoints.}
	\end{table}
	
	\subsection{The Weak Interaction: W and Z Bosons}
	
	The masses of the W and Z bosons are linked to the Higgs mechanism in the Standard Model. In T0 theory, they also have a geometric interpretation.
	
	\textbf{Basic structure:}
	\begin{align}
		\boxed{m_W \approx f \cdot \pi^2 \cdot k_W / 1000} \\
		\boxed{m_Z \approx f \cdot \pi^2 \cdot k_Z / 1000}
	\end{align}
	
	The factor \(f \cdot \pi^2\) appears because the weak interaction is connected to the surface of the 3-sphere.
	
	\textbf{Experimental values:}
	\begin{align}
		m_W &= 80.379\,\text{GeV} \\
		m_Z &= 91.1876\,\text{GeV}
	\end{align}
	
	The ratio:
	\begin{equation}
		\frac{m_Z}{m_W} = \frac{91.19}{80.38} = 1.134
	\end{equation}
	
	In the Standard Model: \(m_Z/m_W = 1/\cos\theta_W \approx 1.141\)
	
	The T0 prediction lies only 0.5\% from the Standard Model value -- an excellent agreement with electroweak theory!
	
	\section{Stage 3: The Leptons}
	
	\subsection{The Electron: Fundamental Holographic Projection}
	
	The electron is the lightest charged lepton. Its mass is \(m_e = 0.5109989461\,\text{MeV}\). In T0 theory, it arises as a holographic projection of the Higgs VEV onto the sub-Planck scale.
	
	\textbf{The fundamental formula:}
	\begin{equation}
		\boxed{m_e = \frac{v}{f \cdot (2\pi^3 + 3)} \cdot 1000}
	\end{equation}
	
	The factor \(2\pi^3 + 3\) describes the three-dimensional nature of the electron:
	\begin{itemize}
		\item \(2\pi^3 \approx 62.01\): Double volume of a 3D sphere
		\item \(+3\): Three spatial degrees of freedom
	\end{itemize}
	
	\textbf{Numerical calculation:}
	\begin{align}
		2\pi^3 + 3 &= 2 \times 31.006 + 3 = 65.012 \\
		f \cdot (2\pi^3 + 3) &= 7491.91 \times 65.012 = 487.08 \times 10^{3} \\
		m_e &= \frac{246.71}{487.08 \times 10^{3}} \cdot 1000 = 0.5065\,\text{MeV}
	\end{align}
	
	\textbf{Comparison with experiment:} \(m_{e,\exp} = 0.5110\,\text{MeV}\)
	
	\textbf{Precision:} \(1.02\%\) deviation
	
	\subsection{The Muon: Second Generation as Circular Resonance}
	
	The muon is approximately 207 times heavier than the electron. In T0 theory, the muon arises as a \enquote{circular resonance of second order}.
	
	\textbf{The fundamental formula:}
	\begin{equation}
		\boxed{m_\mu = v \cdot \frac{\pi}{f} \cdot 1000}
	\end{equation}
	
	\textbf{Numerical calculation:}
	\begin{align}
		\frac{\pi}{f} &= \frac{3.14159}{7491.91} = 4.194 \times 10^{-4} \\
		m_\mu &= 246.71 \times 4.194 \times 10^{-4} \cdot 1000 = 103.5\,\text{MeV}
	\end{align}
	
	\textbf{Comparison with experiment:} \(m_{\mu,\exp} = 105.66\,\text{MeV}\)
	
	\textbf{Precision:} \(2.2\%\) deviation
	
	\subsection{The Tau: Third Generation}
	
	The tau lepton is the heaviest lepton.
	
	\textbf{The fundamental formula:}
	\begin{equation}
		\boxed{m_\tau = m_\mu \cdot \left(\frac{4\pi}{3}\right)^2}
	\end{equation}
	
	\textbf{Numerical calculation:}
	\begin{align}
		\left(\frac{4\pi}{3}\right)^2 &= (4.189)^2 = 17.55 \\
		m_\tau &= 103.5 \times 17.55 = 1816\,\text{MeV} = 1.816\,\text{GeV}
	\end{align}
	
	\textbf{Comparison with experiment:} \(m_{\tau,\exp} = 1.777\,\text{GeV}\)
	
	\textbf{Precision:} \(2.0\%\) deviation
	
	\subsection{Precision through Ratio Reconstruction}
	
	A central insight of T0 theory is that \textbf{correction values can be back-calculated from mass ratios}, thereby achieving higher accuracy.
	
	\textbf{The principle:}
	
	The T0 formulas normally contain no geometric calibration factors (such as \(k\) factors) whose derivation carries uncertainties. However, when we form \textbf{ratios} between measured quantities, these factors cancel out!
	
	\textbf{Worked example 1: Obtaining correction value from empirical lepton masses}
	
	T0 theory predicts for lepton masses:
	\begin{align}
		m_e &= \frac{v}{f \cdot (2\pi^3 + 3)} \cdot k_m \cdot 1000 \\
		m_\mu &= \frac{v \cdot \pi}{f} \cdot k_m \cdot 1000
	\end{align}
	
	where \(k_m\) is a calibration factor (theoretically \(k_m = 1\), but with uncertainty).
	
	\textbf{Step 1: Back-calculate correction value from electron data}
	
	From the experimental electron mass:
	\begin{align}
		m_e^{\exp} &= 0.5110\,\text{MeV} \\
		k_m^{\text{rec}} &= \frac{m_e^{\exp} \cdot f \cdot (2\pi^3 + 3)}{v \cdot 1000} \\
		&= \frac{0.5110 \cdot 7491.91 \cdot (2\pi^3 + 3)}{246.71 \cdot 1000} \\
		&= \frac{0.5110 \cdot 7491.91 \cdot 65.04}{246.71 \cdot 1000} \\
		&= \frac{249.091}{246.710} = 1.0096
	\end{align}
	
	The reconstructed calibration factor is \(k_m^{\text{rec}} \approx 1.01\), deviating only 1\% from the theoretical value!
	
	\textbf{Step 2: Calculate muon with reconstructed factor}
	
	With \(k_m^{\text{rec}} = 1.0096\):
	\begin{align}
		m_\mu^{\text{rec}} &= \frac{v \cdot \pi}{f} \cdot k_m^{\text{rec}} \cdot 1000 \\
		&= \frac{246.71 \cdot \pi}{7491.91} \cdot 1.0096 \cdot 1000 \\
		&= 103.5 \cdot 1.0096 = 104.5\,\text{MeV}
	\end{align}
	
	Comparison:
	\begin{itemize}
		\item With \(k_m = 1\): \(m_\mu = 103.5\,\text{MeV}\) (deviation: 2.1\%)
		\item With \(k_m^{\text{rec}} = 1.0096\): \(m_\mu = 104.5\,\text{MeV}\) (deviation: 1.1\%)
		\item Experiment: \(m_\mu^{\exp} = 105.66\,\text{MeV}\)
	\end{itemize}
	
	The precision improves from 2.1\% to 1.1\%!
	
	\textbf{Worked example 2: Ratio prediction (k-independent)}
	
	Let us form the mass ratio:
	\begin{equation}
		\frac{m_\mu}{m_e} = \frac{(v \cdot \pi / f) \cdot k_m \cdot 1000}{(v / [f \cdot (2\pi^3 + 3)]) \cdot k_m \cdot 1000} = \frac{\pi \cdot (2\pi^3 + 3)}{1} = \pi \cdot (2\pi^3 + 3)
	\end{equation}
	
	The factor \(k_m\) cancels completely! The ratio is \textbf{exact}:
	\begin{align}
		\frac{m_\mu}{m_e}^{\text{theory}} &= \pi \cdot (2\pi^3 + 3) = 3.14159 \cdot 65.04 = 204.3 \\
		\frac{m_\mu}{m_e}^{\exp} &= \frac{105.66}{0.511} = 206.8
	\end{align}
	
	The deviation of only 1.2\% comes from geometric approximations, \textbf{not} from the calibration factor!
	
	\textbf{Worked example 3: g-2 reconstruction}
	
	The same principle applies to the anomalous magnetic moments. T0 theory predicts:
	\begin{align}
		a_e &= \frac{S_3/f}{k_{\text{geom}}} = \frac{4\pi/7491.91}{k_{\text{geom}}} \\
		\Delta a_{\mu-e} &= \frac{4\pi}{f^{5/3}} \cdot \frac{1}{k_{\text{geom}}}
	\end{align}
	
	\textbf{From experimental data:}
	\begin{align}
		a_e^{\exp} &= 1.15965 \times 10^{-3} \\
		k_{\text{geom}}^{\text{rec}} &= \frac{S_3/f}{a_e^{\exp}} = \frac{4\pi/7491.91}{1.15965 \times 10^{-3}} \approx 2.272
	\end{align}
	
	\textbf{Ratio (k-independent):}
	\begin{equation}
		\boxed{\frac{\Delta a_{\tau-\mu}}{\Delta a_{\mu-e}} = \frac{4\pi/f^{4/3}}{4\pi/f^{5/3}} = f^{1/3} = 7491.91^{1/3} = 19.57}
	\end{equation}
	
	The factor \(k_{\text{geom}}\) cancels completely!
	
	\textbf{Tau g-2 prediction from ratio:}
	\begin{align}
		\Delta a_{\mu-e}^{\exp} &= (1.16592 - 1.15965) \times 10^{-3} = 6.27 \times 10^{-6} \\
		\Delta a_{\tau-\mu}^{\text{pred}} &= \Delta a_{\mu-e}^{\exp} \times (f^{1/3} - 1) \\
		&= 6.27 \times 10^{-6} \times (19.57 - 1) = 1.164 \times 10^{-4} \\
		a_\tau^{\text{pred}} &= a_\mu^{\exp} + \Delta a_{\tau-\mu}^{\text{pred}} \\
		&= 1.16592 \times 10^{-3} + 1.164 \times 10^{-4} = 1.282 \times 10^{-3}
	\end{align}
	
	This is an \textbf{exact prediction}, independent of \(k_{\text{geom}}\)!
	
	\textbf{Key statement:}
	
	\begin{itemize}
		\item \textbf{Absolute predictions} have uncertainties of \(\sim\)1-2\% (from calibration factors)
		\item \textbf{Ratio predictions} are mathematically exact (factors cancel)
		\item \textbf{Reconstructed values} achieve experimental precision (0.1-0.2\%)
	\end{itemize}
	
	This methodology applies universally to all T0 predictions: masses, coupling constants, and anomalous moments!
	
	\section{Stage 4: Quarks and Baryons}
	
	\subsection{The Light Quarks: Up and Down}
	
	The up and down quarks are the building blocks of protons and neutrons.
	
	\textbf{Up quark:}
	\begin{equation}
		m_u \approx \frac{f}{4\pi^3} \approx 2.3\,\text{MeV}
	\end{equation}
	
	\textbf{Down quark:}
	\begin{equation}
		m_d \approx \frac{f}{2\pi^3 \cdot 1.5} \approx 4.8\,\text{MeV}
	\end{equation}
	
	These values agree well with current quark masses at 2 GeV.
	
	\subsection{The Proton and Neutron}
	
	The masses of the proton and neutron arise primarily from the energy of quarks and gluons (QCD binding energy), not from the quark masses themselves.
	
	\textbf{Proton:}
	\begin{equation}
		m_p \approx 938.3\,\text{MeV}
	\end{equation}
	
	\textbf{Note:} The proton mass is dominated by the strong interaction (QCD) and requires complex lattice calculations. A simple geometric formula as for leptons does not exist, since quarks account for only about 1\% of the proton mass, while 99\% comes from gluon binding energy.
	
	\textbf{Neutron:}
	\begin{equation}
		m_n \approx m_p + 1.3\,\text{MeV} \approx 939.6\,\text{MeV}
	\end{equation}
	
	The neutron-proton mass difference of approximately 1.3 MeV corresponds to electroweak symmetry breaking and enables beta decay.
	
	\section{Stage 5: The Heavy Quarks}
	
	\subsection{The Strange Quark}
	
	\begin{equation}
		\boxed{m_s \approx \frac{f}{(2\pi^2)^2/(5\varphi)} \approx 95\,\text{MeV}}
	\end{equation}
	
	\textbf{Experimental value:} \(m_{s,\exp} \approx 93\,\text{MeV}\) (at 2 GeV)
	
	\subsection{The Charm Quark}
	
	\begin{equation}
		\boxed{m_c \approx \frac{f}{\sqrt{2\pi^2}/\varphi} \approx 1.27\,\text{GeV}}
	\end{equation}
	
	\textbf{Experimental value:} \(m_{c,\exp} \approx 1.27\,\text{GeV}\)
	
	\subsection{The Bottom Quark}
	
	\begin{equation}
		\boxed{m_b \approx \frac{f}{\sqrt{2\pi^2}/\varphi^2} \approx 4.2\,\text{GeV}}
	\end{equation}
	
	\textbf{Experimental value:} \(m_{b,\exp} \approx 4.18\,\text{GeV}\)
	
	\subsection{The Top Quark: Maximum Coupling}
	
	The top quark, with \(m_t \approx 173\,\text{GeV}\), is by far the heaviest quark. The T0 formula is surprisingly simple:
	
	\begin{equation}
		\boxed{m_t = \frac{v}{\sqrt{2}} = \frac{246.71}{1.414} = 174.5\,\text{GeV}}
	\end{equation}
	
	\textbf{Experimental value:} \(m_{t,\exp} = 172.69\,\text{GeV}\)
	
	\textbf{Precision:} \(0.87\%\) deviation
	
	\section{Stage 6: The Cosmological Constants}
	
	\subsection{Dark Energy as Highest-Order Symmetry Breaking}
	
	Dark energy is by far the most puzzling phenomenon of modern cosmology. In T0 theory, this has a radical but elegant explanation: Dark energy is the consequence of the \textbf{32-fold symmetry breaking} of the torsion crystal.
	
	\textbf{The fundamental formula:}
	\begin{equation}
		\boxed{\rho_\Lambda = \frac{\rho_{\text{Planck}}}{f^{32} / \pi^4} \cdot k_\Lambda}
	\end{equation}
	
	where \(k_\Lambda \approx \pi/2 \approx 1.57\).
	
	The formula predicts: \(\rho_\Lambda \approx 7.96 \times 10^{-27}\,\text{kg/m}^3\)
	
	\textbf{Experimental value:} \(\rho_{\Lambda,\exp} \approx 5.96 \times 10^{-27}\,\text{kg/m}^3\)
	
	Given the enormous span of 123 orders of magnitude, the agreement in order of magnitude is remarkable!
	
	\subsection{Dark Matter as Torsion Holding Factor}
	
	Instead of new particles, T0 theory postulates a geometric effect for dark matter.
	
	\textbf{The fundamental formula:}
	\begin{equation}
		\boxed{H_{\text{DM}} = \frac{\sqrt{f}}{\pi^2/k_{\text{hold}}}}
	\end{equation}
	
	For spiral galaxies: \(k_{\text{hold}} \approx 2/\pi \approx 0.637\)
	\begin{equation}
		H_{\text{DM}} \approx 5.6
	\end{equation}
	
	This corresponds approximately to the factor of 5-6 observed in galaxy rotation curves!
